\section{Related Literature and Positioning}
\label{sec:literature}

Vine copulas are widely used when marginal distributions and dependence must be modelled separately. Pair-copula constructions make high-dimensional dependence tractable while allowing family selection at individual edges \citep{aas_pair-copula_2009,dissmann_selecting_2012}. In portfolio applications, this flexibility is valuable because lower-tail co-movement, asymmetry, and conditional dependence can matter for downside risk \citep{jondeau_copula-garch_2006,sahamkhadam_portfolio_2023}. Most such studies nevertheless transform the fitted distribution into a one-period optimisation problem.

RL portfolio research instead focuses on sequential control. Deep networks can map high-dimensional states to continuous portfolio weights, and recurrent architectures can encode information sequences \citep{jiang_deep_2017,sezer_financial_2020}. Deterministic policy gradients provide a practical baseline for continuous actions \citep{lillicrap_continuous_2015}. However, impressive backtests in this literature are vulnerable to data snooping, unstable hyperparameter selection, and mismatched comparison protocols. Risk sensitivity must also be defined precisely: CVaR is a coherent tail-risk measure \citep{artzner_coherent_1999,rockafellar_optimization_2000}, but a reward penalty is not the same as solving a constrained risk-sensitive control problem.

Our work sits at this intersection. The implementation supplies neural correlation forecasts for every unconditional and conditional edge of a full Student-$t$ D-vine and a two-stage training design for recurrent TD3. The core empirical question remains narrow and testable: conditional on identical long--short exposure and cost constraints, do fully dynamic copula states and synthetic pre-training improve economically meaningful out-of-sample utility over strong non-RL and RL controls? The required controls are specified in \cref{sec:results}.
